\section{Zigzags and Shortening}

\subsection{Shortening and Short Sticklebacks}

If $F=\zigzagTuple{f}{n}$ is a zigzag and 
if for some $i$,
$1 \leq i \leq n-1$,
there is a morphism
$s$ such that $f'_i \circ s = f_{i+1}$ so that the triangle in 
\[
\zigzagShortenable{x}{y}{f}{n}{i}{s}
\]
commutes then the zigzag 
\[
\arraycolsep=7pt
\zigzagShortened{x}{y}{f}{n}{i}{s}
\]
is said to be a \textit{shortening} of $F$ and is said to be \textit{at position i of F}.

\newcommand{\rngTightener}[2]
    {
          \rng{
                \rst{#1_{#2+1}} \circ #1'_{#2}
               }            
    }

\newcommand{\rngTightenerOne}[1]
    {
          \rng{
                \rst{#1_{2}} \circ #1'_{1}
               }            
    }

\newcommand{ \rstTightener}[2]
    {
          \rst{
                #1'_{#2} \circ \rng{#1_{#2}}  
               }
    } 
\begin{notebox}[Change to the above June 24 2026]
Change the above. If $F$ is \emph{not} tight
then need the shortening to be not $f_i \circ s$ as is shown in the diagram but to be \oldt{$f_i \circ \rngTightener{f}{i} \circ s$}  
\newt{$f_i \circ \rng{f'_i} \circ s$}.
\end{notebox}

\begin{observation}
\llabel{shorteningATightZigzagObservation}
If $F$ is tight then in the shortening the
morphism $f_i \circ \rng{f'_i} \circ s$ simplifies to $f_i \circ s$.
\end{observation}
Significantly, the shortening of a tight zigzag is  tight.
It is easy to check this.


Cockett shows that if the category \catcw is a range category
satisfying the addition condition RR.5 then every tight
zigzag of morphisms in \catcw has a unique shortening. 

A stickleback which cannot be shortened is said to be
a \textit{short stickleback}.

\begin{lemma}
\llabel{secondShorteningCharacterisationLemma}
\end{lemma}
If $Z$ is a zigzag of length $2$ and $s$ is a shortening of $Z$
as shown here
\[
\zigzagTwo{x}{y}{z}
\begin{arrows}
\ncarr[15]{y1}{y2} \alabel{s}
\end{arrows}
\]
then
\begin{enumerate}[(i)]
\item 
$\shorten{\tighten{Z}}{s}{1} = \shorten{Z}{s}{1}$
\item 
$\abrst{Z} = \rst{\shorten{Z}{s}{i}}$ 
\item
$\abrng{Z} = \rng{\shorten{Z}{s}{i}}$ 
\end{enumerate}
\begin{proof}
By definition, the shortening of $Z$
is the singleton zigzag
$\tuple{z_1 \circ \rng{z'_1} \circ s}$.

On the otherhand, from  lemma \lref{howToTighteningaZigZagOfLengthTwo}
the tightening $\tighten{Z}$ is the zigzag
$\tuple{ \twoTighteningOne{z},\,\twoTighteningTwo{z},\,\twoTighteningThree{z}}$.
From observation \lref{shorteningATightZigzagObservation}
it then follows that the shortening $\shorten{\tighten{Z}}{s}{i}$
is the singleton zigzag $\tuple{\twoTighteningOne{z} \circ s}$.
\begin{enumerate}[(i)]
\item %{i}
Since $z_1 \circ \rng{\rst{z_2} \circ z'_1} \circ s$ 
simplifies to $z_1 \circ \rng{z'_1} \circ s$ as follows
\begin{align*}
 z_1 \circ \rng{\rst{z_2} \circ z'_1} \circ s
           &= z_1 \circ \rng{\rst{z'_1 \circ s} \circ z'_1} \circ s
           && \mbox{substituting $z'_1 \circ s$ for $z_2$,} \\
           &= z_1 \circ \rng{z'_1 \circ \rst{s}} \circ s 
           && \mbox{by R.4,} \\
           &= z_1 \circ \rng{z'_1} \circ \rst{s} \circ s
           && \mbox{by RR.3,} \\
           &= z_1 \circ \rng{z'_1} \circ s
           && \mbox{by R.1z} 
\end{align*}
then we have
\begin{align*}
\tighten{\shorten{Z}{s}{i}}
            &= \tuple{ z_1 \circ \rng{z'_1} \circ s } \\
            &= \tuple{
                       z_1 \circ \rng{\rst{z_2} \circ z'_1} \circ s
                     } \\
            &= \shorten{\tighten{Z}}{s}{i} &&\mbox{as required.}
\end{align*}  
\item %(ii)
We need to show that 
$\rst{z_1 \circ \rng{z'_1} \circ s}= \rst{z_1 \circ \rngTightenerOne{z}}$.
First we show that 
$\rst{\rng{z'_1} \circ s}=\rst{\rngTightenerOne{z}}$
as follows
\begin{align*}
\rst{\rng{z'_1} \circ s} 
    &= \rst{\rng{z'_1} \circ \rst{s}} && \mbox{wR.4}\\
    &= \rst{ \rng{ z'_1 \circ \rst{s} } }
                                      && \mbox{RR.3} \\
    &= \rst{ \rng{ \rst{z'_1 \circ s} \circ z'_1 } }
                                      && \mbox{R.4} \\
        &= \rst{ \rng{ \rst{z_2} \circ z'_1 } }
                        && \parbox[t]{5cm}{from the initial assumption 
                        $z'_1 \circ s = z_2$.} 
\end{align*}
Now using this we get
\begin{align*}
\rst{z_1 \circ \rng{z'_1} \circ s} 
                &=\rst{z_1 \circ \rst{\rng{z'_1} \circ s}} 
                &&\mbox{by wR.4,}\\
                &=\rst{z_1 \circ \rst{\rngTightenerOne{z}}}
                && \mbox{as shown above,}\\
                &= \rst{z_1 \circ \rngTightenerOne{z}}
                &&\mbox{by wR.4.}
\end{align*}

\item %(iii)
We need to show that 
$\rng{z_1 \circ \rng{z'_1} \circ s} 
     =\rng{\rstTightener{z}{1} \circ z_2}$ 
which we can do as follows
\begin{align*}
\rng{z_1 \circ \rng{z'_1} \circ s} &= \rng{z_1 \circ \rng{z'_1} \circ s}\\
         &= \rng{z'_1 \circ \rng{z_1} \circ s}
            && \mbox{by lemma \lref{fglemma} (vii),}\\
         &= \rng{z'_1 \circ \rst{\rng{z_1}} \circ s}
            && \mbox{by RR.1,}\\
         &= \rng{
                    \rstTightener{z}{1}
                    \circ z'_1 \circ s
                 }
            &&\mbox{by R.4,}\\
         &= \rng{\rstTightener{z}{1} \circ z_2 }
            &&\parbox[t]{5cm}{from the initial assumption  $z'_1 \circ s = z_2$.} 
\end{align*}
\end{enumerate}
\end{proof}

\subsection{Prelude to a category of short zigzags}

In what follows lemma 
\lref{deferredShorteningLemma}
describes the interaction of shortening and tightening 
and is the key to showing that there is a category of short sticklebacks. Before we get to this lemma we have a series of
subordinate lemmas that are required for its proof. 
 
\begin{lemma}
\llabel{shorteningSpecialisationLemma}
In any restriction category \catcw, 
if $f' \leq g'$ and $f \leq g$ then if  
$\nudgeup{1.5cm} % moves the triangle down to make room for arc
\downTriangleShape{x_i}{y_{i-1}}{y_i}
\downToLeft{g'}[0.6]
\downToRight{g}
\leftToRightArcing[a]{s}{30}{30} 
$ commutes 
and $\rst{f}=\rst{f'}$ then
$\downTriangleShape{x_i}{y_{i-1}}{y_i}
\downToLeft{f'}[0.6]
\downToRight{f}
\leftToRightArcing[a]{s}{30}{30} 
$ commutes.
\end{lemma}
\begin{proof}
\begin{align}
f' \circ s 
   &= \rst{f'} \circ g'\circ s 
       && \mbox{because $f' \leq g'$,}\\
   &= \rst{f} \circ g 
   && \mbox{since $\rst{f'}=\rst{f}$ 
              and $g' \circ s = g$, }\\
   &= f
   &&\mbox{since $f \leq g$}
\end{align}
\end{proof}


For convenience, 
whenever the concatenations are defined, we write
$Z \cncdot j$ for $Z \cnc \tuple{j}$ and
$i \dotcnc Z$ for $\tuple{i} \cnc Z$.


$\lefttherefore$

\begin{corollary}
\llabel{shorteningAtightening}
If $F$ is a zigzag and $s$ is a shortening of $F$ at position $i$,then
$s$ is a shortening of 
the tightening of $F$ at position$i$.
\end{corollary}


\begin{lemma}
\llabel{lemmaOfThe6thOfJuly}
\workt{from whiteboard}
If $s$ is a shortening of a zigzag $Z$  of length $2$
then
\begin{enumerate}[(i)]
\item
 if $j$ is a restriction idempotent 
 such that $Z \cncdot j$  is defined
 then $s \circ j$ is a shortening of $Z \cncdot j$ and
 \[
 \shorten{Z \cncdot j}{s \circ j}{1} 
                       = \shorten{Z}{s}{1} \cncdot j,
 \] 
 \item
 if $i$ is a restriction idempotent 
 such that $i \dotcnc Z$  is defined
 then $s$ is a shortening of $i \dotcnc Z$ and
 \[
 \shorten{i \dotcnc Z}{s}{1}
                        = i \dotcnc \shorten{Z}{s}{1}.
 \] 
\end{enumerate}
Note the lack of symmetry between (i) and (ii). 
\end{lemma}
\begin{proof}
\end{proof}

\begin{aside}
If I talk about shortening then the context
needs to be \highlight{RR.5} range category.
\end{aside}

\begin{aside}
I can write  $F= F^{L} \cnc F^{L_r}$

and thus $F\cnc\tuple{g} = (F\cnc\tuple{g})^{L} \cnc (F\cnc\tuple{g})^{L_r}$??

\end{aside}

From this lemma we can see that if a zigzag $F$ shortens then $T(F)$ shortens.

%zigzagIsolationTailTuple{f}{i}{n}
\newcommand{\zigzagIsolationTailTuple}[3]
{
\tuple{\rng{#1_{#2+1}},#1'_{#2+1},...#1'_{#3-1},#1_{#3}} 
}


It is pretty clear that the position of a shortening with a zigzag
can be isolated as described in here
\begin{lemma}
\llabel{isolationOfShortening}
If $F= \zigzagTuple{f}{n}$ is a zigzag and if $i$ indexes a zag
with $F$, i.e. if $1 \leq i \leq n$, then
 
\begin{enumerate}[(i)]
\item if zigzag $L$ of length $i$, zigzag $Z$ of length $2$ and zigzag
$R$ of length $n-i-1$
are the following zigzags
\begin{align*}
L &= \zigzagTuple{f}{i} \\
Z &= \tuple {\rng{f_i}, f'_i, f_{i+1}} \\
R &= \zigzagIsolationTailTuple{f}{i}{n}
\end{align*}
then
\[
F = L \cnc Z \cnc R
\]
\item if $s$ is a shortening of $F$ at position $i$ and then
$s$ is a shortening of $Z$ at position $1$ then
\[
\shorten{F}{s}{i} = L \cnc \shorten{Z}{s}{1} \cnc R
\]
\end{enumerate}
\end{lemma}
\begin{proof}
\end{proof}



This is the key construction used in the proof of the following:
\begin{lemma}
\llabel{TSExtractionLemma}
If $F$ is a zigzag and $s$ is a shortening of $F$ at a position $i$ then
\commentary{conflict: in proof I come up with $s'$ but earlier lemma suggests/proves that $s$ is a shortening of $tighten{F}$. }
\workt{there exists $s'$ which is a shortening of $tighten{F}$???}
\[
\tighten{\shorten{F}{s}{i}} = \shorten{\tighten{F}}{s}{i}
\]
\end{lemma}
\begin{proof}
\commentary{Using abbreviated notation where we write 
               such as $Z_1 \cnc \rst{Z_2}$  for what,
            strictly, should be $Z_1 \cnc \tuple{\rstdot{Z_2}}$.}
We use the construction just given so that we have
\[
F = L \cnc Z \cnc R
\]
where 
$Z$ is of length $2$ and where $L$ and $R$ are zigzags of length 
$\geq 1$, and $s$ is a shortening of $Z$.


We can then use the tightening of $Z$ and lemmas
\lref{howToTightenAThreefoldConcatenationTwo} and
\lref{secondShorteningCharacterisationLemma}.
\workt{sketched out on whiteboard with F for L and G for R}.


\newcommand{\sZ}{\shorten{Z}{s}{1}}
\newcommand{\tL}{\tighten{L}}
\newcommand{\tR}{\tighten{R}}
\newcommand {\leftPartOneOfThree}
{
    \tighten{L 
             \cnc
             \rst{\sZ \cnc \rst{\tR}}
             }
}
\newcommand {\leftPartTwoOfThree}
{
    \rng{\tL} \cnc \sZ \cnc \rst{\tL}
}
\newcommand {\leftPartThreeOfThree}
{
    \tighten{\rng{\rng{\tL} \cnc \sZ}
             \cnc
             R
            }
}


We consider each side of the target identity. For the left hand side we have
\begin{align*}
\tighten{\shorten{F}{s}{i}} 
    &= \tighten{\shorten{L \cnc Z \cnc R}{s}{i}} \\
    &= \tighten{L \cnc \sZ \cnc R} 
       &&\mbox{by lemma \lref{isolationOfShortening},}\\
    &= \mathrlap{
                 \leftPartOneOfThree 
                 \cnc 
                 \leftPartTwoOfThree 
                 \cnc \leftPartThreeOfThree
                 } \\
       &&& \mbox{by lemma \lref{howToTightenAThreefoldConcatenationTwo}}
\end{align*}
where we let
\begin{align*}
l   &= \rng{\tL}           \\
r   &= \rst{\tR}           \\
P_1 &=  \leftPartOneOfThree\\
P_2 &=  \leftPartTwoOfThree\\
P_3 &=  \leftPartThreeOfThree
\end{align*}

\newcommand{\ZZL}{\tighten{Z \cnc r}}
\newcommand{\ZZ} {\tighten{l\cnc Z \cnc r }}
\newcommand{\ZZR}{\tighten{l \cnc Z}}


\newcommand{\rightInteriorPartTwoOfThree}
                 {\shorten{\ZZ}{s \circ r }{i}}

\newcommand{\rightPartOneOfThree}
                 {\tighten{L \cnc \rst{\ZZL}}}
\newcommand{\rightPartTwoOfThree}
                 {\shorten{\ZZ}{s \circ r }{i}}
\newcommand{\rightPartThreeOfThree}
                 {\tighten{\rng{\ZZR} \cnc R}}

Remember:
\begin{align*}
l   &= \rng{\tL}\\
r   &= \rst{\tR}
\end{align*}

For the right hand side let

\begin{align*}
\shorten{\tighten{F}}{s}{i}
   &= \shorten{\tighten{L \cnc Z \cnc R}}
                {s}{i} \\
   &= \mathrlap{
               \shorten{\rightPartOneOfThree
                            \cnc
                           \ZZ 
                            \cnc 
                        \rightPartThreeOfThree
                          }{s}{i} 
               } \\
       &&& \mbox{by lemma \lref{howToTightenAThreefoldConcatenationTwo}. } \\
   &= \mathrlap{
               \rightPartOneOfThree
                    \cnc
                \rightPartTwoOfThree
                    \cnc 
                \rightPartThreeOfThree
               } \\
       &&& \mbox{by a.n.other lemma
                     re:shortening and $\cnc$}\\
       &&& \mbox{and because, by \lref{lemmaOfThe6thOfJuly},}\\
       &&& \mbox{ $s \circ r$ is a shortening of
                    first $l::Z::r$ }\\
       &&& \mbox{ and then  $\ZZ$ }\\
       &&& \mbox{ by corollary \lref{shorteningAtightening}, }\\
    &= Q_1 \cnc Q_2 \cnc Q_3 
\end{align*}
where $Q_1$, $Q_2$, $Q_3$ are given by
\begin{align*}
Q_1 &= \rightPartOneOfThree, \\[0.1cm]
Q_2 &= \rightPartTwoOfThree, \\[0.1cm]
\text{and }Q_3 &= \rightPartThreeOfThree.
\end{align*}
We can complete the proof by showing that $P_1 = Q_1$,
$P_2=Q_2$ and $P_3=Q_3$.
For the first of these ...
\begin{align*}
Q_1 &= \rightPartOneOfThree 
\end{align*}

What seems to be required next in the above is
the (hard fought) 
\lref{secondShorteningCharacterisationLemma}.
In particular since $s \circ \rst{r}$ is a shortening
of $Z \conc \rst{r}$, since $r$ is an idempotent then
\[
\rst{\ZZ} 
   = \shorten{Z \cnc r}{s \circ r}{1}
\]


For the second ..
\begin{align*}
Q_2 &= \rightPartTwoOfThree 
\end{align*}
For the third...
\begin{align*}
Q_3 &= \rightPartThreeOfThree 
\end{align*}

\end{proof}
\begin{aside}
The following lemma suggests that the normal form for a zigzag is the
shortening of the tightening.
\end{aside}
\begin{lemma}
\llabel{deferredShorteningLemma}
If $Z_1$, $Z_2$ are tight zigzags in
\newt{an RR.5 range category} such that the composition $Z_1 :: Z_2$ is defined then
\begin{enumerate}[(i)]
\item 
$shorten(tighten(shorten(Z_1) :: Z_2))= shorten(tighten(Z_1 :: Z_2)).$
\item
$shorten(tighten(Z_1 :: shorten(Z_2)))= shorten(tighten(Z_1 :: Z_2)).$
\end{enumerate}
\end{lemma}
\begin{proof}
Since in both (i) and (ii) the innermost shortening is a combination 
of individual shortenings and since
it follows by lemma \lref{associativitySublemma} that 
each of these individual shortenings can be moved outside of the tightening.
\end{proof}

\subsection{The category of short sticklebacks}

If \catcw is an RR.5 range catgeory then
 there is a category of short sticklebacks over \catc, which we denote
 $\Stk$.
Composition of short sticklebacks is defined by
\[
F \circ G = shorten(T(F::G)) 
\]
Composition, so defined, is associative because:
\begin{align*}
F \circ (G \circ H)
  &= shorten(T(F :: shorten(T(G ::H)) && \mbox{definition $\circ$,} \\
  &= shorten (T(F::T(G::H))) && \mbox{by lemma \ref{deferredShorteningLemma} (i),} \\
  &= shorten (T(F::G::H))    && \mbox{\parbox[t]{5cm}{see proof of associativity of composition of tight zigzags, section \ref{categorytightzigzags}.}}
\end{align*}  
and because, likewise, 
using lemma \ref{deferredShorteningLemma} (ii),
we have also that $(F \circ G) \circ H=shorten (T(F::G::H))$.

\subsection{Sequential Shortening}
\begin{lemma}
\llabel{sequentialShorteningLemma}
If $F=\zigzagTuple{f}{n}:x_1 \morph y_n$ 
is a tight stickleback in \highlight{an RR.5 range category} \catcw
and if $F$ shortens to a zigzag $\tuple{f}$ for some
morphism $f:x_1 \morph y_n$ in \catcw then  there
is a sequence of morphisms $s_1,...s_{n-1}$,
$s_i: y_i \morph y_n$ in \catcw
such that all diagrams in 
\commentary{Improve the diagramming at a later date}
\[
%\arraycolsep=12pt
\zigzagSequentialShortening{x}{y}{f}{n}{x}{g}{h}{s}
\]

commute i.e. such that
\[
f_1 \circ s_1 =f
\]
and \foreachi[n-2], 
\[
f'_i \circ s_i =  f_{i+1} \circ s_{i+1}
\] 
and
\[
f'_{n-1} \circ s_{n-1} = f_n 
\]
\end{lemma}
\begin{proof}
For any series of shortenings of $F$ to a single morphism one of this set must shorten the triangle the final zag in F i.e. 
there must be a morphism $s_{n-1}:y_{n-1} \morph y_n$
such that the triangle 
$\nudgeup{1.5cm}\downTriangleShape{x_n}{y_{n-1}}{y_n}
\downToLeft{g'}[0.6]
\downToRight{g}
\leftToRightArcing[a]{s}{30}{30} 
$ 
commutes. 
Since shortenings can be applied in any order this shortening can be applied first and the resulting zigzag 
$\tuple{f_1,f'_1,...f_{n-2}, f_{n-1} \circ s_{n-1}}$
must then shorten to the morphism $f$. 
Repeating the argument we get a morphism $s_{n-2}$
such that 
$
\begin{array}{c  c c}
&\Ob{y}{n}&               \\[0.5cm]
\Obpp{y}{n} && \Obp{y}{n} \\[0.5cm]
& \Obp{x}{n}
\end{array}
\begin{arrows}
\ncarr{ynpp}{yn}\alabel{s_{n-2}}
\ncarr{ynp}{yn}\blabel{s_{n-1}}
\ncarr{xnp}{ynpp}\alabel{f'_{n-2}}
\ncarr{xnp}{ynp}\blabel{f_{n-1}}
\end{arrows}
$
commutes and shorten $F$ to a zigzag
$\tuple{f_1,f'_1,...f_{n-3}, f_{n-2} \circ s_{n-2}}$
By iteration we arrive
at morphisms $s_1,...s_{n-1}$ which shorten $F$ to a singleton (and therefore short) zigzag $\tuple{f_1 \circ s_1}$. But since the shortening of a tight zigzag is unique 
we must have $f_1 \circ s_1=f$.
\end{proof}
